% !TeX spellcheck = de_DE

Der entwickelte Ansatz basiert auf einer zyklischen Abfolge aus Bewertung und Überarbeitung.
Ausgangspunkt ist ein bereits erzeugter LS+-Text, dieser wird zunächst anhand von Metriken analysiert.
Die ermittelten Werte werden anschließend mit zuvor festgelegten Anforderungen verglichen.

Sind die Anforderungen erfüllt, gilt der Text als qualitativ ausreichend und der Prozess wird beendet.
Gibt es jedoch eine Differenz zwischen den erreichten Metriken und den Soll-Metriken, werden die erkannten Schwachstellen in Form von Feedback aufbereitet und für eine erneute Generierung des Textes verwendet.
Das Large Language Model erhält dabei gezielte Hinweise zu den Bereichen, die verbessert werden sollen.

Im nächsten Schritt werden die Texte anschließend erneut bewertet.
Dieser Vorgang wiederholt sich, bis entweder die definierten Qualitätsanforderungen erfüllt sind oder die maximale Anzahl an Iterationen erreicht wird.

Durch diesen Ansatz soll die Textoptimierung nicht ausschließlich dem Sprachmodell überlassen werden.
Stattdessen erfolgt die Überarbeitung auf Basis von berechenbaren Metriken, die die Verbesserungsschleife kontrolliert lenken.